%% kapa: exam=2 round=34 year=2023 period=1 subject=appraisal_practice subject_ko=감정평가실무 seq=1 number=1 points=40
%% source: Q-Net 감정평가사 기출문제 (cst00309) artlSeq=5212168 file=raw/1008/5212168/1교시 23년도 감정평가사 2차시험_감정평가실무.pdf
%% conversion: kordoc 4.24.51 -> tools/gen2.py (auto); review=proofread-gpt-sol; flags=none
\begin{longquestion}{【문제 1】}
감정평가사 甲은 (주)A자산운용으로부터 현황 부동산(이하 ‘대상부동산’) 및 최유효이용 상정 부동산(이하 ‘최유효이용 부동산’)에 대해 감정평가를 의뢰받았다. 제시된 자료를 활용하여 다음 물음에 답하시오. (40점) \par
\medskip\noindent 물음 1) 공시지가기준법 및 원가법으로 대상부동산을 감정평가하시오. (5점) \par
\medskip\noindent 물음 2) 거래사례비교법으로 대상부동산을 감정평가하시오. (5점) \par
\medskip\noindent 물음 3) 수익환원법(직접환원법)으로 대상부동산을 감정평가하시오. (5점) \par
\medskip\noindent 물음 4) 물음 1)～물음 3)의 시산가액을 기준으로 대상부동산의 감정평가액을 결정하시오. (5점) \par
\medskip\noindent 물음 5) 할인현금흐름분석법을 적용하여 최유효이용 부동산을 감정평가하시오. (15점) \par
\medskip\noindent 물음 6) 최유효이용에 미달하는 부분의 가치를 구하시오. (5점) \par
< 자료 1 > 대상부동산 및 최유효이용 부동산 개요 \par
1. 기준가치: 시장가치 \par
2. 기준시점: 2023.07.15 \par
3. 대상부동산 개요 \par
1) 토지 \par
{\small\centering
\resizebox{\ifdim\width>\linewidth\linewidth\else\width\fi}{!}{%
\begin{tabular}{|c|c|c|c|c|c|c|c|}
\hline
\makecell{소재지} & \makecell{지목} & \makecell{면적\\(m²)} & \makecell{용도\\지역} & \makecell{이용\\상황} & \makecell{도로\\접면} & \makecell{형상\\지세} & \makecell{주위환경} \\
\hline
\makecell{K구 S동\\100} & \makecell{대} & \makecell{800} & \makecell{일반상업} & \makecell{업무용} & \makecell{광대\\소각} & \makecell{세장형\\평지} & \makecell{업무지대} \\
\hline
\end{tabular}}\par}

2) 건물 \par
{\small\centering
\resizebox{\ifdim\width>\linewidth\linewidth\else\width\fi}{!}{%
\begin{tabular}{|c|c|c|c|c|c|c|}
\hline
\makecell{소재지} & \makecell{구조} & \makecell{층수} & \makecell{연면적\\(m²)} & \makecell{용도} & \makecell{급수} & \makecell{비고} \\
\hline
\makecell{K구 S동\\100} & \makecell{철근\\콘크리트조} & \makecell{지하2층/\\지상5층} & \makecell{2,700} & \makecell{업무용} & \makecell{3} & \makecell{사용승인일:\\1985.07.10} \\
\hline
\end{tabular}}\par}

4. 최유효이용 부동산 개요 \par
1) 토지 \par
{\small\centering
\resizebox{\ifdim\width>\linewidth\linewidth\else\width\fi}{!}{%
\begin{tabular}{|c|c|c|c|c|c|c|c|}
\hline
\makecell{소재지} & \makecell{지목} & \makecell{면적\\(m²)} & \makecell{용도\\지역} & \makecell{이용\\상황} & \makecell{도로\\접면} & \makecell{형상\\지세} & \makecell{주위환경} \\
\hline
\makecell{K구 S동\\100} & \makecell{대} & \makecell{800} & \makecell{일반상업} & \makecell{업무용} & \makecell{광대\\소각} & \makecell{세장형\\평지} & \makecell{업무지대} \\
\hline
\end{tabular}}\par}

2) 건물 \par
{\small\centering
\resizebox{\ifdim\width>\linewidth\linewidth\else\width\fi}{!}{%
\begin{tabular}{|c|c|c|c|c|c|c|}
\hline
\makecell{소재지} & \makecell{구조} & \makecell{층수} & \makecell{연면적\\(m²)} & \makecell{용도} & \makecell{급수} & \makecell{비고} \\
\hline
\makecell{K구 S동\\100} & \makecell{철골철근\\콘크리트조} & \makecell{지하5층/\\지상18층} & \makecell{9,600} & \makecell{업무용} & \makecell{1} & \makecell{사용승인 예정일:\\2025.07.15} \\
\hline
\end{tabular}}\par}

< 자료 2 > 표준지공시지가 \par
(공시기준일: 2023.01.01) \par
{\small\centering
\resizebox{\ifdim\width>\linewidth\linewidth\else\width\fi}{!}{%
\begin{tabular}{|c|c|c|c|c|c|c|c|c|}
\hline
\makecell{기호} & \makecell{소재지} & \makecell{지목} & \makecell{면적\\(m²)} & \makecell{이용\\상황} & \makecell{용도지역} & \makecell{도로\\교통} & \makecell{형상\\지세} & \makecell{공시지가\\(원/m²)} \\
\hline
\makecell{1} & \makecell{K구 S동\\90} & \makecell{대} & \makecell{900} & \makecell{상업용} & \makecell{일반상업} & \makecell{중로\\한면} & \makecell{정방형\\평지} & \makecell{54,000,000} \\
\hline
\makecell{2} & \makecell{K구 S동\\120} & \makecell{대} & \makecell{850} & \makecell{업무용} & \makecell{일반상업\\3종일주} & \makecell{광대\\한면} & \makecell{부정형\\평지} & \makecell{80,000,000} \\
\hline
\makecell{3} & \makecell{K구 S동\\140} & \makecell{대} & \makecell{1,200} & \makecell{업무용} & \makecell{근린상업} & \makecell{광대\\세각} & \makecell{정방형\\평지} & \makecell{90,000,000} \\
\hline
\makecell{4} & \makecell{K구 S동\\160} & \makecell{대} & \makecell{1,000} & \makecell{업무용} & \makecell{중심상업} & \makecell{광대\\소각} & \makecell{정방형\\평지} & \makecell{100,000,000} \\
\hline
\makecell{5} & \makecell{K구 S동\\190} & \makecell{대} & \makecell{900} & \makecell{업무용} & \makecell{일반상업} & \makecell{광대\\소각} & \makecell{정방형\\평지} & \makecell{90,000,000} \\
\hline
\makecell{6} & \makecell{S구 R동\\150} & \makecell{대} & \makecell{1,250} & \makecell{업무용} & \makecell{중심상업} & \makecell{광대\\세각} & \makecell{정방형\\평지} & \makecell{95,000,000} \\
\hline
\makecell{7} & \makecell{S구 R동\\200} & \makecell{대} & \makecell{700} & \makecell{상업용} & \makecell{일반상업} & \makecell{중로\\각지} & \makecell{정방형\\평지} & \makecell{68,000,000} \\
\hline
\end{tabular}}\par}

※ 1) S구 R동은 K구 S동과 동일수급권 내 유사지역에 소재함 \par
2) 기호 1, 7은 노선 상가지대, 기호 2～6은 업무지대에 소재함 \par
3) 일반상업과 근린상업, 일반상업과 중심상업, 근린상업과 중심상업은 지역격차가 있음 \par
4) 기호 2의 3종일주 면적은 30\%임 \par
< 자료 3 > 인근지역 평가사례 및 거래사례 \par
1. 평가사례(평가목적: 일반거래) \par
{\small\centering
\resizebox{\ifdim\width>\linewidth\linewidth\else\width\fi}{!}{%
\begin{tabular}{|c|c|c|c|c|c|c|c|c|c|}
\hline
\makecell{기호} & \makecell{소재지} & \makecell{지목} & \makecell{면적\\(m²)} & \makecell{용도지역} & \makecell{이용\\상황} & \makecell{도로\\접면} & \makecell{형상\\지세} & \makecell{기준\\시점} & \makecell{토지단가\\(원/m²)} \\
\hline
\makecell{가} & \makecell{K구 S동\\80} & \makecell{대} & \makecell{770} & \makecell{일반상업} & \makecell{상업용} & \makecell{중로한면} & \makecell{정방형\\평지} & \makecell{2022.02.01} & \makecell{55,500,000} \\
\hline
\makecell{나} & \makecell{K구 S동\\115} & \makecell{대} & \makecell{3,000} & \makecell{일반상업\\3종일주} & \makecell{업무용} & \makecell{광대\\한면} & \makecell{세장형\\평지} & \makecell{2023.03.03} & \makecell{85,000,000} \\
\hline
\makecell{다} & \makecell{K구 S동\\185} & \makecell{대} & \makecell{2,000} & \makecell{일반상업} & \makecell{상업용} & \makecell{광대\\한면} & \makecell{가장형\\평지} & \makecell{2020.06.01} & \makecell{80,000,000} \\
\hline
\makecell{라} & \makecell{K구 S동\\200} & \makecell{대} & \makecell{850} & \makecell{일반상업} & \makecell{업무용} & \makecell{광대\\한면} & \makecell{세장형\\평지} & \makecell{2023.07.01} & \makecell{90,000,000} \\
\hline
\makecell{마} & \makecell{K구 S동\\250} & \makecell{대} & \makecell{1,360} & \makecell{일반상업} & \makecell{업무용} & \makecell{중로한면} & \makecell{가장형\\평지} & \makecell{2021.01.01} & \makecell{52,000,000} \\
\hline
\end{tabular}}\par}

※ 기호 나 3종일주에 속하는 면적은 미미하여 일반상업으로 판단함 \par
2. 거래사례 \par
1) 거래사례 \#1 \par
- 소재지: K구 S동 70 \par
- 총 거래가격: 45,000,000,000원 \par
- 거래시점: 2023.04.01 \par
- 토지: 일반상업, 상업용, 750m², 중로한면, 세장형, 평지 \par
- 건물 \par
{\small\centering
\resizebox{\ifdim\width>\linewidth\linewidth\else\width\fi}{!}{%
\begin{tabular}{|c|c|c|c|c|}
\hline
\makecell{구조} & \makecell{급수} & \makecell{연면적\\(m²)} & \makecell{사용승인일} & \makecell{부대설비 내역} \\
\hline
\makecell{철근콘크리트조} & \makecell{2} & \makecell{3,000} & \makecell{2016.02.20} & \makecell{대상건물과 유사함} \\
\hline
\end{tabular}}\par}

- 기타사항: 노선 상가지대에 위치하며, 정상 거래사례임 \par
2) 거래사례 \#2 \par
- 소재지: K구 S동 380 \par
- 총 거래가격: 71,400,000,000원 \par
- 거래시점: 2023.02.01 \par
- 토지: 일반상업, 업무용, 840m², 광대세각, 세장형, 평지 \par
- 건물 \par
{\small\centering
\resizebox{\ifdim\width>\linewidth\linewidth\else\width\fi}{!}{%
\begin{tabular}{|c|c|c|c|c|}
\hline
\makecell{구조} & \makecell{급수} & \makecell{연면적\\(m²)} & \makecell{사용승인일} & \makecell{부대설비 내역} \\
\hline
\makecell{철근콘크리트조} & \makecell{3} & \makecell{2,500} & \makecell{1984.01.25} & \makecell{대상건물과 유사함} \\
\hline
\end{tabular}}\par}

- 기타사항: 업무지대에 소재하며, 정상 거래사례임 \par
3) 거래사례 \#3 \par
- 소재지: K구 S동 355 \par
- 총 거래가격: 48,750,000,000원 \par
- 거래시점: 2023.05.01 \par
- 토지: 일반상업, 업무용, 650m², 광대소각, 부정형, 평지 \par
- 건물 \par
{\small\centering
\resizebox{\ifdim\width>\linewidth\linewidth\else\width\fi}{!}{%
\begin{tabular}{|c|c|c|c|c|}
\hline
\makecell{구조} & \makecell{급수} & \makecell{연면적\\(m²)} & \makecell{사용승인일} & \makecell{부대설비 내역} \\
\hline
\makecell{철근콘크리트조} & \makecell{4} & \makecell{2,000} & \makecell{1979.02.19} & \makecell{대상건물과 유사함} \\
\hline
\end{tabular}}\par}

- 기타사항: 업무지대에 위치하며, 건물 노후화로 총 거래가격에 건부감가가 포함되어 있는 사례임 \par
4) 거래사례 \#4 \par
- 소재지: K구 S동 60 \par
- 총 거래가격: 47,740,000,000원 \par
- 거래시점: 2023.04.01 \par
- 토지: 일반상업, 업무용, 770m², 중로한면, 가장형, 평지 \par
- 건물 \par
{\small\centering
\resizebox{\ifdim\width>\linewidth\linewidth\else\width\fi}{!}{%
\begin{tabular}{|c|c|c|c|c|}
\hline
\makecell{구조} & \makecell{급수} & \makecell{연면적\\(m²)} & \makecell{사용승인일} & \makecell{부대설비 내역} \\
\hline
\makecell{철근콘크리트조} & \makecell{4} & \makecell{2,500} & \makecell{1983.03.29} & \makecell{대상건물과 유사함} \\
\hline
\end{tabular}}\par}

- 기타사항: 노선 상가지대에 위치하며, 정상 거래사례임 \par
< 자료 4 > 시점수정 관련 자료 \par
1. 지가변동률(S시 K구) \par
{\small\centering
\resizebox{\ifdim\width>\linewidth\linewidth\else\width\fi}{!}{%
\begin{tabular}{|c|c|c|}
\hline
\makecell{구분} & \makecell{주거지역(\%)} & \makecell{상업지역(\%)} \\
\hline
\makecell{2020.06.01～2023.07.15(누계)} & \makecell{15.345} & \makecell{16.565} \\
\hline
\makecell{2021.01.01～2023.07.15(누계)} & \makecell{11.050} & \makecell{12.115} \\
\hline
\makecell{2022.02.01～2023.07.15(누계)} & \makecell{7.585} & \makecell{7.885} \\
\hline
\makecell{2023.01.01～2023.06.30(누계)} & \makecell{1.270} & \makecell{1.295} \\
\hline
\makecell{2023.02.01～2023.07.15(누계)} & \makecell{1.150} & \makecell{1.165} \\
\hline
\makecell{2023.03.03～2023.07.15(누계)} & \makecell{1.020} & \makecell{1.035} \\
\hline
\makecell{2023.04.01～2023.07.15(누계)} & \makecell{1.050} & \makecell{1.070} \\
\hline
\makecell{2023.05.01～2023.07.15(누계)} & \makecell{1.010} & \makecell{1.020} \\
\hline
\makecell{2023.06.01～2023.06.30} & \makecell{0.100} & \makecell{0.150} \\
\hline
\end{tabular}}\par}

※ 2023년 7월 이후 지가변동률은 미고시 되었음 \par
2. 오피스빌딩 자본수익률(S시 K구) \par
{\small\centering
\resizebox{\ifdim\width>\linewidth\linewidth\else\width\fi}{!}{%
\begin{tabular}{|c|c|c|c|c|}
\hline
\makecell{구분} & \makecell{2022년 3분기} & \makecell{2022년 4분기} & \makecell{2023년 1분기} & \makecell{2023년 2분기} \\
\hline
\makecell{자본수익률(\%)} & \makecell{0.22} & \makecell{0.26} & \makecell{0.20} & \makecell{0.24} \\
\hline
\end{tabular}}\par}

3. 건축비지수는 변동이 없다고 가정함 \par
< 자료 5 > 지역요인 관련 자료 \par
1. 대상부동산과 동일한 구(區)에 소재하는 표준지공시지가, 평가사례 및 거래사례는 지역요인이 동일함 \par
2. 다른 구(區)에 소재하는 표준지공시지가, 평가사례 및 거래사례는 지역요인의 비교가 필요함 \par
< 자료 6 > 개별요인 관련 자료 \par
1. 가로조건(각지인 경우 가로조건에서 반영함) \par
{\small\centering
\resizebox{\ifdim\width>\linewidth\linewidth\else\width\fi}{!}{%
\begin{tabular}{|c|c|c|c|c|c|c|c|}
\hline
\makecell{구분} & \makecell{광대한면} & \makecell{광대소각} & \makecell{광대세각} & \makecell{중로한면} & \makecell{중로각지} & \makecell{소로한면} & \makecell{소로각지} \\
\hline
\makecell{광대한면} & \makecell{1.00} & \makecell{1.10} & \makecell{1.06} & \makecell{0.96} & \makecell{0.98} & \makecell{0.86} & \makecell{0.90} \\
\hline
\makecell{광대소각} & \makecell{0.93} & \makecell{1.00} & \makecell{0.97} & \makecell{0.88} & \makecell{0.92} & \makecell{0.79} & \makecell{0.83} \\
\hline
\makecell{광대세각} & \makecell{0.96} & \makecell{1.05} & \makecell{1.00} & \makecell{0.91} & \makecell{0.95} & \makecell{0.82} & \makecell{0.86} \\
\hline
\makecell{중로한면} & \makecell{1.06} & \makecell{1.16} & \makecell{1.12} & \makecell{1.00} & \makecell{1.05} & \makecell{0.90} & \makecell{0.95} \\
\hline
\makecell{중로각지} & \makecell{1.02} & \makecell{1.11} & \makecell{1.07} & \makecell{0.97} & \makecell{1.00} & \makecell{0.87} & \makecell{0.91} \\
\hline
\makecell{소로한면} & \makecell{1.19} & \makecell{1.29} & \makecell{1.25} & \makecell{1.13} & \makecell{1.17} & \makecell{1.00} & \makecell{1.05} \\
\hline
\makecell{소로각지} & \makecell{1.13} & \makecell{1.23} & \makecell{1.19} & \makecell{1.08} & \makecell{1.12} & \makecell{0.97} & \makecell{1.00} \\
\hline
\end{tabular}}\par}

2. 획지조건 \par
{\small\centering
\resizebox{\ifdim\width>\linewidth\linewidth\else\width\fi}{!}{%
\begin{tabular}{|c|c|c|c|c|c|}
\hline
\makecell{구분} & \makecell{정방형} & \makecell{가장형} & \makecell{세장형} & \makecell{사다리형} & \makecell{부정형} \\
\hline
\makecell{정방형} & \makecell{1.00} & \makecell{1.00} & \makecell{0.97} & \makecell{0.93} & \makecell{0.91} \\
\hline
\makecell{가장형} & \makecell{1.00} & \makecell{1.00} & \makecell{0.97} & \makecell{0.93} & \makecell{0.91} \\
\hline
\makecell{세장형} & \makecell{1.02} & \makecell{1.02} & \makecell{1.00} & \makecell{0.95} & \makecell{0.93} \\
\hline
\makecell{사다리형} & \makecell{1.06} & \makecell{1.06} & \makecell{1.04} & \makecell{1.00} & \makecell{0.96} \\
\hline
\makecell{부정형} & \makecell{1.10} & \makecell{1.10} & \makecell{1.08} & \makecell{1.04} & \makecell{1.00} \\
\hline
\end{tabular}}\par}

3. 상기 제시된 조건 이외의 다른 조건은 동일함 \par
< 자료 7 > 토지, 건물 일괄 거래사례비교법 적용 개별요인 관련 자료 \par
1. 개별요인 \par
1) 대상물건 / 거래사례 \#1 \par
{\small\centering
\resizebox{\ifdim\width>\linewidth\linewidth\else\width\fi}{!}{%
\begin{tabular}{|c|c|c|c|}
\hline
\makecell{구분} & \makecell{입지적 특성} & \makecell{기능적 특성} & \makecell{물리적 특성} \\
\hline
\makecell{대상물건} & \makecell{115} & \makecell{100} & \makecell{95} \\
\hline
\makecell{거래사례 \#1} & \makecell{100} & \makecell{100} & \makecell{100} \\
\hline
\end{tabular}}\par}

2) 대상물건 / 거래사례 \#2 \par
{\small\centering
\resizebox{\ifdim\width>\linewidth\linewidth\else\width\fi}{!}{%
\begin{tabular}{|c|c|c|c|}
\hline
\makecell{구분} & \makecell{입지적 특성} & \makecell{기능적 특성} & \makecell{물리적 특성} \\
\hline
\makecell{대상물건} & \makecell{104} & \makecell{100} & \makecell{100} \\
\hline
\makecell{거래사례 \#2} & \makecell{100} & \makecell{100} & \makecell{100} \\
\hline
\end{tabular}}\par}

3) 대상물건 / 거래사례 \#3 \par
{\small\centering
\resizebox{\ifdim\width>\linewidth\linewidth\else\width\fi}{!}{%
\begin{tabular}{|c|c|c|c|}
\hline
\makecell{구분} & \makecell{입지적 특성} & \makecell{기능적 특성} & \makecell{물리적 특성} \\
\hline
\makecell{대상물건} & \makecell{100} & \makecell{100} & \makecell{105} \\
\hline
\makecell{거래사례 \#3} & \makecell{100} & \makecell{100} & \makecell{100} \\
\hline
\end{tabular}}\par}

4) 대상물건 / 거래사례 \#4 \par
{\small\centering
\resizebox{\ifdim\width>\linewidth\linewidth\else\width\fi}{!}{%
\begin{tabular}{|c|c|c|c|}
\hline
\makecell{구분} & \makecell{입지적 특성} & \makecell{기능적 특성} & \makecell{물리적 특성} \\
\hline
\makecell{대상물건} & \makecell{115} & \makecell{100} & \makecell{105} \\
\hline
\makecell{거래사례 \#4} & \makecell{100} & \makecell{100} & \makecell{100} \\
\hline
\end{tabular}}\par}

2. 토지, 건물 일괄 개별요인은 상승식으로 계산하되 소수점 셋째자리에서 반올림하여 둘째자리까지 표기하고, 제시된 특성 이외의 특성은 동일함 \par
< 자료 8 > 재조달원가 및 감가수정 관련자료 \par
1. 표준단가 \par
{\small\centering
\resizebox{\ifdim\width>\linewidth\linewidth\else\width\fi}{!}{%
\begin{tabular}{|c|c|c|c|c|}
\hline
\makecell{용도} & \makecell{구조} & \makecell{급수} & \makecell{표준단가(원/m²)} & \makecell{내용연수(년)} \\
\hline
\multirow{2}{*}{\makecell{업무시설}} & \makecell{철근콘크리트조\\(5층～15층 이하)} & \makecell{1} & \makecell{2,000,000} & \makecell{50} \\
\cline{2-5}
 & \makecell{철골철근콘크리트조\\(10층～20층 이하)} & \makecell{1} & \makecell{2,300,000} & \makecell{55} \\
\hline
\multirow{2}{*}{\makecell{업무시설}} & \makecell{철근콘크리트조\\(5층～15층 이하)} & \makecell{2} & \makecell{1,700,000} & \makecell{50} \\
\cline{2-5}
 & \makecell{철골철근콘크리트조\\(10층～20층 이하)} & \makecell{2} & \makecell{2,000,000} & \makecell{55} \\
\hline
\makecell{업무시설} & \makecell{철근콘크리트조\\(5층～15층 이하)} & \makecell{3} & \makecell{1,500,000} & \makecell{50} \\
\hline
\makecell{업무시설} & \makecell{철근콘크리트조\\(5층～15층 이하)} & \makecell{4} & \makecell{1,300,000} & \makecell{50} \\
\hline
\makecell{업무시설} & \makecell{철근콘크리트조\\(5층～15층 이하)} & \makecell{5} & \makecell{1,200,000} & \makecell{50} \\
\hline
\end{tabular}}\par}

2. 부대설비 보정단가는 상기 표준단가에 포함되었음 \par
3. 건물 잔가율은 0\%임 \par
4. 건물의 감가수정은 정액법(만년감가)을 적용함 \par
5. 최유효이용 건물의 건축비는 표준단가를 적용함 \par
< 자료 9 > 대상부동산 및 최유효이용 부동산의 시장임대료 관련 자료 \par
1. 대상부동산 관련 시장임대료 \par
{\small\centering
\resizebox{\ifdim\width>\linewidth\linewidth\else\width\fi}{!}{%
\begin{tabular}{|c|c|c|c|c|}
\hline
\makecell{구분} & \makecell{임대사례 1} & \makecell{임대사례 2} & \makecell{임대사례 3} & \makecell{대상부동산} \\
\hline
\makecell{소재지} & \makecell{S시\\K구 S동} & \makecell{S시\\K구 S동} & \makecell{S시\\K구 S동} & \makecell{S시\\K구 S동 100} \\
\hline
\makecell{건물명} & \makecell{가 빌딩} & \makecell{나 빌딩} & \makecell{다 빌딩} & \makecell{○○빌딩} \\
\hline
\makecell{층(지상/지하)} & \makecell{17F/B3} & \makecell{5F/B2} & \makecell{5F/B2} & \makecell{5F/B2} \\
\hline
\makecell{구조} & \makecell{철골철근\\콘크리트} & \makecell{철근\\콘크리트} & \makecell{철근\\콘크리트} & \makecell{철근\\콘크리트} \\
\hline
\makecell{건물연면적(m²)} & \makecell{10,000} & \makecell{2,800} & \makecell{2,600} & \makecell{2,700} \\
\hline
\makecell{토지면적(m²)} & \makecell{900} & \makecell{850} & \makecell{750} & \makecell{800} \\
\hline
\makecell{용적률(\%)} & \makecell{800} & \makecell{230} & \makecell{210} & \makecell{220} \\
\hline
\makecell{용도지역} & \makecell{일반상업} & \makecell{일반상업} & \makecell{일반상업} & \makecell{일반상업} \\
\hline
\makecell{사용승인(년)} & \makecell{2022} & \makecell{1986} & \makecell{1987} & \makecell{1985} \\
\hline
\makecell{오피스빌딩\\하위시장} & \makecell{B북부} & \makecell{B북부} & \makecell{B남부} & \makecell{B북부} \\
\hline
\makecell{보증금(원/m²)} & \makecell{300,000} & \makecell{270,000} & \makecell{260,000} & \makecell{250,000} \\
\hline
\makecell{월임대료(원/m²)} & \makecell{30,000} & \makecell{27,000} & \makecell{26,000} & \makecell{25,000} \\
\hline
\makecell{월관리비(원/m²)} & \makecell{15,000} & \makecell{12,000} & \makecell{11,000} & \makecell{12,000} \\
\hline
\makecell{비고} & \makecell{시장임대료} & \makecell{시장임대료} & \makecell{시장임대료} & \makecell{계약임대료} \\
\hline
\end{tabular}}\par}

※ 1) 오피스빌딩 하위시장이 다른 경우 그 격차는 < 자료 10 >을 기준으로 보정해야 함 \par
2) 시장임대료 및 계약임대료는 기준시점에서 조사된 임대료이며, 연면적을 기준으로 함 \par
2. 최유효이용 부동산 관련 시장임대료 \par
{\small\centering
\resizebox{\ifdim\width>\linewidth\linewidth\else\width\fi}{!}{%
\begin{tabular}{|c|c|c|c|c|}
\hline
\makecell{구분} & \makecell{임대사례 4} & \makecell{임대사례 5} & \makecell{임대사례 6} & \makecell{최유효이용\\부동산} \\
\hline
\makecell{소재지} & \makecell{S시\\K구 S동} & \makecell{S시\\K구 S동} & \makecell{S시\\K구 S동} & \makecell{S시\\K구 S동 100} \\
\hline
\makecell{건물명} & \makecell{라 빌딩} & \makecell{마 빌딩} & \makecell{바 빌딩} & \makecell{○○빌딩} \\
\hline
\makecell{층(지상/지하)} & \makecell{19F/B5} & \makecell{17F/B5} & \makecell{18F/B5} & \makecell{18F/B5} \\
\hline
\makecell{구조} & \makecell{철골철근\\콘크리트} & \makecell{철골철근\\콘크리트} & \makecell{철골철근\\콘크리트} & \makecell{철골철근\\콘크리트} \\
\hline
\makecell{건물연면적(m²)} & \makecell{10,000} & \makecell{9,000} & \makecell{12,000} & \makecell{9,600} \\
\hline
\makecell{토지면적(m²)} & \makecell{950} & \makecell{850} & \makecell{1,000} & \makecell{800} \\
\hline
\makecell{용적률(\%)} & \makecell{799} & \makecell{780} & \makecell{800} & \makecell{800} \\
\hline
\makecell{용도지역} & \makecell{일반상업} & \makecell{일반상업} & \makecell{일반상업} & \makecell{일반상업} \\
\hline
\makecell{사용승인(년)\\(사용승인예정)(년)} & \makecell{2021} & \makecell{2022} & \makecell{2018} & \makecell{(2025)} \\
\hline
\makecell{오피스빌딩 하위시장} & \makecell{B남부} & \makecell{B북부} & \makecell{B북부} & \makecell{B북부} \\
\hline
\makecell{보증금(원/m²)} & \makecell{270,000} & \makecell{300,000} & \makecell{290,000} & \makecell{-} \\
\hline
\makecell{월임대료(원/m²)} & \makecell{27,000} & \makecell{30,000} & \makecell{29,000} & \makecell{-} \\
\hline
\makecell{월관리비(원/m²)} & \makecell{14,000} & \makecell{15,000} & \makecell{15,000} & \makecell{-} \\
\hline
\makecell{비고} & \makecell{시장임대료} & \makecell{시장임대료} & \makecell{시장임대료} & \makecell{-} \\
\hline
\end{tabular}}\par}

※ 1) 오피스빌딩 하위시장이 다른 경우 그 격차는 < 자료 10 >을 기준으로 보정해야 함 \par
2) 시장임대료는 기준시점에서 조사된 최유효이용 부동산 관련 시장임대료이며, 연면적을 기준으로 함 \par
< 자료 10 > 수익환원법 적용 관련 자료 \par
1. 오피스빌딩 하위시장의 격차를 보정하는 자료는 조사가 불가능함 \par
2. 렌트프리(Rent Free)는 없음 \par
3. 임대사례와 대상부동산 및 최유효이용 부동산의 개별요인은 동일함 \par
4. 직접환원법 및 할인현금흐름분석법에 적용할 조건은 다음과 같음 \par
1) 환원율은 투자결합법으로 산출한 결과 연 4.5\%임 \par
2) 보증금 운용수익률은 연 3\%, 공실손실상당액 및 대손충당금은 가능총수익의 5\%, 보증금ㆍ연간 임대료ㆍ연간 관리비 상승율은 연 2\%, 할인율은 시장에서 발표된 부동산투자수익률 및 물가상승률을 고려할 때 연 5\%임 \par
3) 직접환원법에 적용할 운영경비는 연간 관리비 중 70\%, 할인현금흐름분석법에 적용할 운영경비는 연간 관리비 중 60\%임 \par
4) 복귀가액은 내부추계법을 적용하며, 재매도비용은 2\%임 \par
5) 대상부동산이 속한 B북부 오피스빌딩 하위시장은 최유효이용 부동산이 사용승인 된 후 인근지역에 GTX-A 노선 역이 신설 예정이고, C그룹 본사의 오피스빌딩이 신축되는 등 개발호재로 인해 복귀가액 산정을 위한 최종환원율은 할인율에서 0.5\%p를 공제하여 적용하는 것이 타당한 것으로 조사됨 \par
6) 할인현금흐름분석법은 최유효이용 부동산을 5년 보유 후 6년차에 매도하는 것으로 가정함 \par
7) 할인율 및 기별 현재가치 계수는 다음과 같음 \par
{\small\centering
\resizebox{\ifdim\width>\linewidth\linewidth\else\width\fi}{!}{%
\begin{tabular}{|c|c|c|c|c|c|c|}
\hline
\makecell{구분} & \makecell{1기} & \makecell{2기} & \makecell{3기} & \makecell{4기} & \makecell{5기} & \makecell{6기} \\
\hline
\makecell{4.0\%} & \makecell{0.961} & \makecell{0.924} & \makecell{0.889} & \makecell{0.854} & \makecell{0.821} & \makecell{0.790} \\
\hline
\makecell{4.5\%} & \makecell{0.956} & \makecell{0.915} & \makecell{0.876} & \makecell{0.838} & \makecell{0.802} & \makecell{0.767} \\
\hline
\makecell{5.0\%} & \makecell{0.952} & \makecell{0.907} & \makecell{0.863} & \makecell{0.822} & \makecell{0.783} & \makecell{0.746} \\
\hline
\makecell{5.5\%} & \makecell{0.947} & \makecell{0.898} & \makecell{0.851} & \makecell{0.807} & \makecell{0.765} & \makecell{0.725} \\
\hline
\end{tabular}}\par}

< 자료 11 > 기타 자료 \par
1. 표준지공시지가 및 사례 선정 시 선정사유를 반드시 기재 \par
2. 개별요인은 상승식으로 산출하되, 소수점 셋째자리에서 반올림하여 둘째자리까지 표기 \par
3. 그 밖의 요인 보정치는 “표준지 기준 산정방식”을 적용 \par
4. 각 감정평가방법의 시산가액은 천만원 단위 이하에서 절사하여 억원 단위로 표기 \par
5. 최유효이용에 미달하는 거래사례의 총 거래가격은 토지 면적을 기준으로 거래되어 시점수정은 < 자료 4 > 시점수정 관련 자료를 적용 \par
6. 할인현금흐름분석법의 현금흐름표는 1～6년 차를 모두 기재 \par
7. 대상부동산의 감정평가액은 공시지가기준법 및 원가법 시산가액의 40\%, 거래사례비교법 시산가액의 30\%, 수익환원법 시산가액의 30\%로 적용하여 조정 및 결정 \par
8. 대상부동산 건물의 철거비는 고려하지 않음 \par
\end{longquestion}
